\section{Introduction}
\label{sec:intro}

Quantitative finance has invested heavily in prediction and comparatively little in the
step that follows it. A large body of work now forecasts returns, directions or
volatilities with machine learning, and reports its success in statistical terms:
accuracy, error, or the significance of a coefficient \citep{gu20, zhang22}. Far less
work specifies what an investor should then hold. That second step is not a
formality. A forecast becomes an investment only after it has been reconciled with a
covariance structure, a risk budget, position limits and the cost of trading, and the
translation can destroy the value of a forecast as easily as preserve it. This paper
refers to the missing step as the \emph{signal--allocation gap}.

\subsection{The Signal--Allocation Gap}

The gap is visible in how results are reported. Systematic reviews of fusion-based
prediction find that statistical accuracy remains the dominant success criterion, with
economic value reported inconsistently or not at all \citep{zhang22}. Where allocation
is addressed, it is often through heuristics rather than optimisation: the IntPort
framework of \citet{sami25}, for instance, converts forecasts into weights by
statistical averaging, without a covariance structure or an explicit risk objective.
Reinforcement learning approaches optimise a policy directly and so avoid the
translation problem, but the majority are formulated for a single instrument
\citep{saijai26, jiang25}, which leaves the portfolio-level question open.

This matters because the properties an investor cares about are portfolio properties.
Diversification, drawdown and the stability of holdings are not visible at the level of
an individual forecast, and \citet{demiguel09} showed that optimisation which ignores
estimation error can underperform naive equal weighting out of sample. Any framework
that claims to bridge prediction and allocation must therefore be judged on portfolio
outcomes under realistic frictions, not on the quality of its forecasts.

\subsection{Regime Dependence}

Financial markets do not behave homogeneously through time. Since \citet{hamilton89},
regime-switching models have been used to represent discrete shifts in the statistical
behaviour of financial series, and \citet{angbekaert02} demonstrated that accounting for
such shifts materially changes optimal international allocations. \citet{kritzman12}
made the practical case that regime awareness alters what a dynamic strategy should do,
not merely how it should be described. More recent work continues in this direction,
combining regime identification with allocation rules \citep{li25, luo26}.

Two consequences follow for a framework of the kind proposed here. The first is that a
predictive signal's reliability is conditional: an indicator that is informative in a
quiet trending market need not be informative when the market is disturbed. The second
is that the appropriate amount of risk to carry is itself state-dependent, which is the
premise behind volatility-managed strategies \citep{moreira17, hood25} and behind risk
budgeting more generally \citep{asimit25}. A framework that detects regimes but acts on
them in only one of these two ways will leave part of its own premise untested.

\subsection{What This Paper Does}

This paper specifies a three-stage decision-theoretic framework and tests it end to
end. Latent market states are identified by a Gaussian mixture model fitted on an
expanding window. Two composite signals---one technical, one derived from price and
volume dynamics---are fused with regime-conditional weights and converted to expected
excess returns using information-ratio scaling. The fused signal then enters a
constrained quadratic program with Ledoit--Wolf covariance shrinkage
\citep{ledoit04} and an explicit turnover penalty, and the resulting sleeve is scaled by
a regime-dependent risk budget.

The framework is evaluated on ten liquid exchange-traded funds spanning equities, fixed
income, credit, gold and real estate over eighteen years, under a walk-forward protocol
with transaction costs. The evaluation is designed to answer not only whether the
framework works but which of its parts is responsible, and it reports the answer in
both directions. The regime risk budget is shown to be the effective component: at
identical average exposure, timing the de-risking by regime improves both the Sharpe
ratio and the maximum drawdown. The signal fusion layer is shown not to be: the
information coefficients of both composites are statistically indistinguishable from
zero or significantly negative, and the regime-conditional weighting hypothesis is
rejected by the data that motivated it.

Reporting the second finding is a deliberate choice. Component-level nulls are rarely
published, which is one reason architectures accumulate layers whose contribution has
never been isolated. \citet{harvey15} and \citet{valls26} document how easily an
unreported specification search inflates apparent performance in this setting, and
\citet{demiguel09} showed the practical cost of trusting an unvalidated component. An
ablation that locates the source of a result is more useful to a practitioner than an
aggregate number that does not.

\subsection{Contributions}

\begin{itemize}

\item A fully specified regime-conditional allocation framework, given in enough detail
to be reproduced from the text alone: exact feature definitions, the mixture model and
its label-anchoring requirement, the optimisation problem with all constraints, and the
walk-forward protocol.

\item A two-step separation of the allocation decision into relative composition and
total risk. The paper shows that expressing the regime risk budget as a fraction of the
portfolio's own volatility, rather than as an absolute volatility target, is what keeps
the regime layer active: absolute targets calibrated for concentrated equity portfolios
bind in one rebalance out of 224 on a diversified universe and render the layer inert.

\item An ablation design that isolates component contributions at matched risk. Because
scaling a portfolio toward cash changes both its return and its volatility, comparing a
de-risked strategy with a fully invested benchmark confounds risk timing with risk
level. Holding average exposure fixed separates the two, and shows that the regime
layer earns its place through timing rather than through simply holding less.

\item A reported null on regime-conditional signal fusion, together with the
information coefficients that produce it, measured overall and within each state.

\item A documented diagnosis of a failure mode in signal-to-allocation conversion:
passing an unscaled indicator into a mean-variance objective allows the ratio of units
between the linear and quadratic terms, rather than the information in the signal, to
determine the solution, producing a degenerate single-asset portfolio.

\end{itemize}

\subsection{Organisation}

Section~2 reviews portfolio construction, regime-aware allocation and the treatment of
textual sentiment, and positions this study against them. Section~3 specifies the
framework. Section~4 reports the results, including robustness, ablation and the
negative findings. Section~5 discusses implications and limitations, and Section~6
concludes.
